%% appendix_missions.tex — The Cassini--Huygens Mission and Dragonfly
%% Extracted from main.tex
\chapter{The Cassini--Huygens Mission and Dragonfly}
\label{app:missions}

The \cassini{}--\huygens{} mission (2004--2017) provided the observational foundation of this thesis through 127 targeted \titan{} flybys, the \huygens{} atmospheric probe descent and landing, and a systematic multi-instrument survey of \titan{}'s surface, atmosphere, and interior.  SAR mapping covered $\sim$67\% of the surface; VIMS near-infrared spectroscopy $\sim$50\%; CIRS thermal mapping the full globe at $\sim$10\textdegree{} resolution.  The mission's key unresolved gap is the absence of a quantitative, spatially-resolved, multi-epoch habitability analysis integrating all data modalities --- which this thesis provides.

NASA's \dragonfly{} rotorcraft \citep{Lorenz2018Dragonfly}, due to arrive in 2034, will traverse $\sim$175\,km from the Belet dune field to Selk crater over a 2.7-year mission.  Chapter~\ref{ch:results} provides a quantitative evaluation of that site selection relative to globally optimal alternatives.
